\chapter{Serving Agents: FastAPI, Streaming and Concurrency}
\label{ch:serving}

\begin{quote}\itshape
This is the chapter the job advertisement was actually about.
\end{quote}

\noindent
Everything so far has been an agent. This chapter is about the service around it:
one process holding hundreds of concurrent conversations, streaming tokens to
browsers, surviving client disconnections, respecting a provider's rate limit,
and shutting down without losing work. It is where Part~I's \code{asyncio}
material and Part~III's persistence material meet, and it is the point at which
a demo becomes a system.

\chapterstub{
\textbf{Argument.} Because state lives in Postgres rather than in the process,
pods are stateless and scale horizontally --- but that same property creates the
one hazard the architecture does not solve for you: two concurrent runs on one
\code{thread\_id}. That needs a lock or a per-thread queue, and it is the thing
most implementations miss.

\textbf{Sections.} (13.1) Why the agent is built once in \code{lifespan} and not
per request. (13.2) Dependency injection in FastAPI, and the shape that stays
testable. (13.3) Streaming to a client: server-sent events versus WebSockets,
with the event format worth committing to. (13.4) Client disconnection ---
cancel the run, or finish it in the background to preserve the checkpoint? Both,
with the cases for each. (13.5) Backpressure: a slow consumer and an unbounded
queue is a memory leak with a delay fuse. (13.6) Rate limiting: a process-wide
semaphore, jittered backoff, and model fallback. (13.7) The three timeout levels
in a request context. (13.8) Blocking tools inside an async service, and the
thread-pool budget. (13.9) Multiple workers: shared checkpointer, stateless pods,
and enforcing one active run per thread. (13.10) Long-running work: queues,
background workers, cron, and why it does not belong in a request.

\textbf{Listings.} The complete service --- \code{lifespan}, chat endpoint, SSE
generator; disconnection handled both ways; a bounded queue between producer and
SSE consumer; a per-thread advisory lock in Postgres; a load test at a few
hundred concurrent conversations.

\textbf{Diagrams.} \code{svc-architecture} --- the deployed system: pods,
Postgres, provider, client. \code{svc-sse-flow} --- sequence diagram of a
streamed turn including a disconnect. \code{svc-thread-lock} --- two concurrent
requests on one thread, with and without the lock.

\textbf{Measurement.} Concurrent conversations sustained by one process against
p50/p95 latency, with a mocked provider so the number measures the service rather
than the model. Then the same with one blocking call introduced, which is the
most persuasive graph in the book.

\textbf{Mini-project.} Stage 10: the FastAPI service --- SSE, disconnect policy,
rate limiting, per-thread locking, load test.

\textbf{Source topics covered.} Part~VI 6.1--6.12.
}
